\subsection{The interaction
Hamiltonian}\label{proof:higher_order_obstructions:the-interaction-hamiltonian}

Let \(L = \operatorname{arctanh} \rho\). At \(\beta  = 1\), an exact
Hamiltonian up to an additive constant is

\[
H_{n,\rho}(x)=-\log\sum_{a\in\mathbb F_2^m}
\cosh\left(L\sum_z x_zc_a(z)\right). \tag{H2}
\]

This follows by summing the product-noise likelihoods over \(A,B\). The
conditional-probability bound implies
\(|H(x)-H(x^{(i)})| \le  2 \operatorname{arctanh} \rho\). Thus
single-spin energy changes can be arbitrarily small and entropy density
arbitrarily close to \(\log  2\), yet AA fails.

The energy oscillation is at most \(\log  \cosh (Ln) \le  Ln\): the sum
of coshes lies between n and \(n \cosh (Ln)\). Thus this obstruction
also respects an \(O(n)\) total-energy budget. It relies on global
higher-order interactions, not a superextensive energy range.

The energy is invariant under coordinatewise multiplication by every
affine character codeword. Its Fourier coefficients can therefore be
nonzero only on sets S with \(|S|\) even and \(\sum_{z\in S}z=0\). The
smallest nonempty such set has size four. There are no one-, two-, or
three-body Fourier terms. This is a higher-order global interaction, not
a four-locality claim: higher orders are generally present.

Consequently uniform marginals, identity covariance, three-wise
independence, near-maximal entropy density, a uniform lower bound on
single-site conditional randomness, and small one-spin energy changes
are jointly insufficient for influence AA. These are genuine negatives
furnished by H1.

They do not refute Dobrushin contraction, spectral independence under
every pinning, or bounded interaction order: those stronger hypotheses
have not been established for this example.
